\begin{afterword}
\summaryhead

The preface introduces a collection of blog posts that the author wrote over two years. It begins by listing five mistakes. The largest was not seeing that learning this way of thinking is mainly a matter of practice, not theory. The second was writing about big, abstract problems instead of everyday life. The third was stressing belief over action. The fourth was poor organization, which Rob Bensinger's reordering of the posts corrects. The fifth was open contempt for ideas the author thought stupid. The author would now write more courteously, while still thinking that mockery can help.

The preface then states the goal: to teach a valuable way of thinking that schools do not teach. It is tied to science and to the willingness to say ``Oops'' and give up on what is not working, in daily life as in the laboratory. Despite the mistakes, the posts appeared to help many people, because it is all, in the end, ``one thing.'' The author concludes that the book is better than nothing.

\responsehead

The preface is candid. It opens the book with a list of the book's largest faults, and the list reads as meant. The difficulty is what follows from it.

By the preface's own account, the book corrects one of the five mistakes, the fourth: the posts have been reordered. The first three concern what the posts teach and how: theory where practice was needed, distant and abstract examples, belief over action. The preface does not say that the reordering, or the ``bit'' of rewriting, addressed any of them. The repair of the mistake about practice is assigned to the Center for Applied Rationality, outside the book. So the reader is told at the start that the text ahead has these faults, and is then given the text.

The case for reading it anyway is thin. The posts ``appeared to help a surprising number of people a surprising amount''; no number, example or source is given. The reply to having ``picked all the wrong examples'' is that ``it is all, in the end, one thing,'' which is asserted twice and not argued.

The goal is described loosely. The way of thinking is said to be ``not taught systematically at all,'' without support. What the preface says it contains is the experimental method and the willingness to say ``Oops'' and give up on what is not working. The rest is left to the book.

On the fifth mistake the preface is divided. It regrets the contempt and says the author would now write more courteously, ``I think.'' But it still holds that ``derisive laughter by comedians can help people wake up from the dream,'' and offers no evidence for that.

In short: a candid list of the book's largest faults, of which the book corrects only the ordering, followed by a case for reading it that rests on the author's impression that it helped.
\end{afterword}
